% Page 030 — Alban pendant la guerre de Cent Ans
% Transcription fidèle de la page fournie.
% Les graphies et formulations du document ont été conservées.

\begin{center}
    {\Large\itshape Alban pendant la guerre de Cent Ans}
\end{center}

\vspace{1.5em}

\noindent
Comment le trésorier du comte d'Armagnac séjournant en ce mois de mars 1371 à Alban,
se tira-t-il de ces « vilains draps » ? L'envoyé des consuls d'Albi retourna-t-il
bredouille ? Peut être heureux de n'avoir pas été molesté par les soudards de ce
grand féodal ? On ne sait.

\medskip

\noindent
Ce qui est sûr, c'est que la guerre n'a apporté que des ruines. Alban s'est tiré
d'affaire sans, apparemment trop de mal. Mais en cet an 1371, la famine sévit et
bien des gens meurent de faim, et la guerre n'est pas finie. En 1380, les consuls
albanais écrivent à ceux d'Albi pour les renseigner sur des mouvements de troupes.
D'Alban, Bernat Domergue, en juin 1381, cours à Albi, on n'a plus le temps d'écrire.
Ça presse. Les Albigeois sont prévenue par notre concitoyen qu'un certain Bénézet,
Capitaine de routiers vient d'occuper Pousthomy, Montfranc, Trébas. En 1382, on
trouve ce brigand confortablement installé à Saint Jean de Janes. Il ne l'évacue
que contre rançon. Le château de Jaufin est pris et évacué en 1384 dans les mêmes
onéreuses conditions. Les dégâts sont tels dans la région d'Alban que sur
19 paroisses englobées dans l'archiprêtré de Nougaret (nous dirions doyenné
aujourd'hui), 7 paroisses sont exonérées d'impôts. Notre Dame d'Alban ne paie que
15 sols aux collecteurs pontificaux. Le Truel, Carménel, Ainestères, Saint Pierre
d'Alban, Massuyès, Ambialet, Villeneuve ne paient rien. C'est la dèche généralisée.
